\documentclass[10pt]{article}

\usepackage[T1]{fontenc}
\usepackage{lmodern}
\usepackage[margin=0.62in,top=0.52in,bottom=0.55in]{geometry}
\usepackage{amsmath,amssymb}
\usepackage{array,booktabs,tabularx}
\usepackage{graphicx}
\usepackage{microtype}
\usepackage{xcolor}
\usepackage[hidelinks]{hyperref}

\definecolor{MSSNavy}{HTML}{12304A}
\definecolor{MSSOrange}{HTML}{C65D16}
\definecolor{MSSGreen}{HTML}{2F855A}
\definecolor{MSSGray}{HTML}{627D98}
\definecolor{MSSPale}{HTML}{F4F7FA}

% Generated by Code/scripts/build_extension_proposal_assets.py.
% Do not edit numerical values by hand.
\newcommand{\ExtensionLiabilityMortgageTopOne}{+4.50}
\newcommand{\ExtensionLiabilityMortgageNextForty}{+23.63}
\newcommand{\ExtensionLiabilityMortgageBottomFifty}{+13.85}
\newcommand{\ExtensionLiabilityConsumerTopOne}{0.00}
\newcommand{\ExtensionLiabilityConsumerNextForty}{+2.92}
\newcommand{\ExtensionLiabilityConsumerBottomFifty}{+4.04}
\newcommand{\ExtensionSavingMortgageTopOne}{-0.43}
\newcommand{\ExtensionSavingMortgageNextForty}{-3.09}
\newcommand{\ExtensionSavingMortgageBottomFifty}{-2.28}
\newcommand{\ExtensionSavingConsumerTopOne}{-0.01}
\newcommand{\ExtensionSavingConsumerNextForty}{-0.73}
\newcommand{\ExtensionSavingConsumerBottomFifty}{-0.88}
\newcommand{\ExtensionNetMortgageTopOne}{+8.69}
\newcommand{\ExtensionNetMortgageNextForty}{-19.02}
\newcommand{\ExtensionNetMortgageBottomFifty}{-13.73}
\newcommand{\ExtensionNetConsumerTopOne}{+2.78}
\newcommand{\ExtensionNetConsumerNextForty}{-3.46}
\newcommand{\ExtensionNetConsumerBottomFifty}{-4.29}
\newcommand{\ExtensionBottomNinetyNineMortgageBoom}{-6.36}
\newcommand{\ExtensionBottomNinetyNineMortgagePost}{-0.16}
\newcommand{\ExtensionBottomNinetyNineConsumerBoom}{-1.76}
\newcommand{\ExtensionBottomNinetyNineConsumerPost}{-1.44}


\newcommand{\compactheading}[1]{%
  \vspace{0.32em}\noindent{\color{MSSOrange}\bfseries\large #1}\par\vspace{0.16em}}
\newcommand{\resulttag}[1]{\textcolor{MSSGreen}{\textbf{#1}}}
\newcolumntype{L}[1]{>{\raggedright\arraybackslash}p{#1}}

\setlength{\parindent}{0pt}
\setlength{\parskip}{0.32em}
\setlength{\tabcolsep}{4pt}
\renewcommand{\arraystretch}{1.12}

\begin{document}

{\color{MSSOrange}\bfseries\small PROPOSED EXTENSION}\hfill
{\color{MSSGray}\small 15 August 2026}\par
{\color{MSSNavy}\bfseries\LARGE What Drives the Saving-Rate Divergence?}\par
{\color{MSSGray}\small Debt composition and ultimate financing across the wealth distribution}
\vspace{0.25em}\hrule height 1.1pt\vspace{0.35em}

\compactheading{Starting point: Figure 6 asks ``why?''}

Mian, Straub, and Sufi show that net saving rates rise steeply with wealth and
that the profile becomes more unequal after 1982. Figure 6 compares the
1963--1982 and 1983--2019 averages; our exact-method reconstruction and its
trailing five-year heatmap add the missing within-period timing. They show that
the divergence is not a uniform secular shift: saving stays persistently high
at the top, while the largest negative movements below the top are concentrated
around the mid-2000s and the Great Recession. This motivates a mechanism question that
the two long-period profiles cannot answer on their own.

The paper already establishes three essential facts. First, Figure 1 documents
the rise of financial intermediation before 2008. Second, Table A2 shows that
greater borrowing is the largest accounting component of the decline in saving
below the top 1 percent. Third, Figure 8 unveils a widening net-lender position
at the top against borrowing lower in the wealth distribution. These results
connect rich saving to household debt in the aggregate, but they do not show
which debt contracts produce the saving-rate change or whether the same
mechanism operates before and after 2008.

\vspace{0.15em}
\noindent\fcolorbox{MSSOrange}{MSSPale}{%
\begin{minipage}{0.955\textwidth}
\textbf{Research question.} Which debt instruments account for the decline in
net saving rates below the top 1 percent, who ultimately finances each
instrument after intermediary claims are unveiled, and how much of the
pre-2008 divergence can be attributed to exogenous mortgage-credit expansion?
\end{minipage}}

\compactheading{Measurement bridge}

For wealth group $g$, debt category $k$, and signed liability balance
$W^{L,k}_{gt}<0$, define the write-down-adjusted active-saving contribution

$$
  \Theta^{L,k}_{gt}
  =W^{L,k}_{gt}-(1+\pi^{L,k}_{gt})W^{L,k}_{g,t-1},
  \qquad
  s^{L,k}_{gt}=\frac{\Theta^{L,k}_{gt}}{Z^d_{gt}},
$$

where $\pi^{L,k}_{gt}$ is the debt valuation factor and $Z^d_{gt}$ is personal
plus attributable corporate disposable income. Negative $\Theta^{L,k}_{gt}$
means active borrowing reduced measured saving. Summing $s^{L,k}_{gt}$ over
debt categories, and then adding financial-asset and real-estate saving, must
reproduce the exact Figure 6 rate. On the lender side, the same category is
traced through the paper's operator $\Omega^k_t=(I-Q_t)^{-1}B^k_t$ before
subtracting group liabilities. This creates one auditable bridge from
\emph{which group borrowed which instrument} to \emph{which group ultimately
financed it}.

\compactheading{Relation to the paper and its broader mechanism}

The proposal turns the paper's accounting link into a debt-instrument-specific
empirical design. It complements \emph{Indebted Demand}, which provides a
general-equilibrium mechanism connecting inequality, household debt, weak
demand, and low interest rates, and the mortgage-credit literature of Mian and
Sufi. It does not presume that greater financialization caused the saving-rate
divergence; instead, it asks which balance-sheet conduit is quantitatively
relevant before selecting a causal shock.

\vfill
{\color{MSSGray}\scriptsize
Source: Mian, Straub, and Sufi, ``The Saving Glut of the Rich,'' July 2025,
Sections 3.4--4.2, Figures 1, 6, and 8, and Table A2; Mian, Straub, and Sufi,
``Indebted Demand,'' 2021. The heatmap referenced above uses the paper's exact
annual rate definition on the February 2021 input vintage through 2016.}

\clearpage
{\color{MSSOrange}\bfseries\small PRELIMINARY EVIDENCE}\hfill
{\color{MSSGray}\small Same-vintage descriptive accounting}\par
{\color{MSSNavy}\bfseries\LARGE Mortgages Select the Pre-2008 Mechanism}\par
{\color{MSSGray}\small February 2021 inputs, full Leontief unveiling, 1963--2016}
\vspace{0.25em}\hrule height 1.1pt\vspace{0.35em}

The current implementation separates home mortgages from aggregate consumer
credit for four wealth groups, applies the paper's group-specific write-down
factors, and unveils each category from saved direct cells. No old seven-round
output enters the calculation. The table reports selected groups to expose the
top-lender/middle-borrower contrast; all entries are percentage points of
aggregate national income.

\begin{tabularx}{\textwidth}{@{}X L{1.05in}rrr@{}}
\toprule
\textbf{Object and period} & \textbf{Debt category} &
\textbf{Top 1\%} & \textbf{Next 40\%} & \textbf{Bottom 50\%} \\
\midrule
\multicolumn{5}{@{}l}{\textit{Change in liability stock, 1982--2007}} \\
& Home mortgages
& $\ExtensionLiabilityMortgageTopOne$
& $\ExtensionLiabilityMortgageNextForty$
& $\ExtensionLiabilityMortgageBottomFifty$ \\
& Consumer credit
& $\ExtensionLiabilityConsumerTopOne$
& $\ExtensionLiabilityConsumerNextForty$
& $\ExtensionLiabilityConsumerBottomFifty$ \\
\addlinespace
\multicolumn{5}{@{}l}{\textit{Average annual contribution to active saving, 1998--2007}} \\
& Home mortgages
& $\ExtensionSavingMortgageTopOne$
& $\ExtensionSavingMortgageNextForty$
& $\ExtensionSavingMortgageBottomFifty$ \\
& Consumer credit
& $\ExtensionSavingConsumerTopOne$
& $\ExtensionSavingConsumerNextForty$
& $\ExtensionSavingConsumerBottomFifty$ \\
\addlinespace
\multicolumn{5}{@{}l}{\textit{Change in unveiled net position, 1982--2007}} \\
& Home mortgages
& $\ExtensionNetMortgageTopOne$
& $\ExtensionNetMortgageNextForty$
& $\ExtensionNetMortgageBottomFifty$ \\
& Consumer credit
& $\ExtensionNetConsumerTopOne$
& $\ExtensionNetConsumerNextForty$
& $\ExtensionNetConsumerBottomFifty$ \\
\bottomrule
\end{tabularx}

\vspace{0.2em}
\resulttag{Motivating result.}
Mortgages dominate the 1982--2007 debt build-up, the 1998--2007 debt-related
saving drag, and the widening lender--borrower split. The top 1 percent was
already a net mortgage lender in 1982; its position widened rather than
switching sign. The regime changes after 2008: for the bottom 99 percent, the
average mortgage contribution moves from
$\ExtensionBottomNinetyNineMortgageBoom$ in 1998--2007 to
$\ExtensionBottomNinetyNineMortgagePost$ in 2008--2016, while consumer credit
moves only from $\ExtensionBottomNinetyNineConsumerBoom$ to
$\ExtensionBottomNinetyNineConsumerPost$. Thus mortgages select the primary
pre-crisis mechanism, whereas post-crisis consumer credit requires a separate
composition analysis.

\compactheading{Proposed research design}

\textbf{1. Close the measurement identity.} Re-express mortgage and
consumer-credit contributions with each group's exact Figure 6 disposable-income
denominator, retain the percentile bins, and verify that debt categories sum
to total debt saving and that all asset classes recover the complete saving
rate. This distinguishes a debt contribution from an explanation of the full
rate.

\textbf{2. Identify borrower incidence and contract type.} Use the Survey of
Consumer Finances to split mortgages and home-equity loans from revolving,
vehicle, student, installment, and other debt. Report participation and
conditional balances by wealth and age, and repeat the analysis using
gross-asset, focal-debt-excluded, or lagged rank so that debt does not mechanically
determine the household's wealth group.

\textbf{3. Estimate mechanisms by regime.} For 1982--2007, combine a validated
policy or lender-exposure shock to mortgage credit with borrower outcomes and
the category-specific unveiling map. For the post-2008 period, first determine
whether student, vehicle, or revolving credit accounts for the residual before
choosing a separate identification strategy. A single aggregate time-series
regression would confound credit with house prices, interest rates,
demographics, regulation, and inequality.

\compactheading{Novelty and evidentiary boundary}

Debt composition by wealth and the pre-crisis mortgage boom are not themselves
new facts. The contribution is the intersection: connecting category-specific
borrower incidence to the exact Figure 6 saving-rate divergence and to Figure
8 ultimate ownership. The preliminary results establish accounting incidence,
not that top saving caused loan origination, mortgage credit caused the full
saving decline, or borrowing reduced welfare. Those claims remain outcomes for
the proposed causal stage.

\vfill
{\color{MSSGray}\scriptsize\raggedright
Generated from the verified debt-composition headline and period-summary
tables.\par
Empirical producer: \texttt{build\_debt\_composition\_extension.py}. Displayed
values: \texttt{build\_extension\_proposal\_assets.py}. The February 2021
authors' package is unchanged.}

\end{document}
